\documentclass[10pt]{article}
\usepackage[utf8]{inputenc}
\usepackage[T1]{fontenc}
\usepackage[margin=0.95in]{geometry}
\usepackage{amsmath,amssymb,amsthm}
\usepackage{booktabs}
\usepackage{array}
\usepackage{url}
\usepackage{microtype}
\usepackage{xcolor}
\usepackage{listings}
\usepackage[hidelinks]{hyperref}
\makeatletter\g@addto@macro{\UrlBreaks}{\do\_\do\.\do\/}\makeatother
\emergencystretch=3em
\tolerance=2000
\lstset{basicstyle=\ttfamily\footnotesize,columns=fullflexible,keepspaces=true,breaklines=true,
  frame=single,framesep=3pt,xleftmargin=4pt,xrightmargin=4pt,upquote=true,aboveskip=4pt,belowskip=2pt}
\newcommand{\tier}[1]{\textsf{[#1]}}
\newcommand{\tA}{\tier{A}}
\newcommand{\tL}{\tier{L}}
\newcommand{\tC}{\tier{C}}
\newcommand{\tX}{\tier{X}}
\newcommand{\lean}[1]{\texttt{#1}}
\newcommand{\id}[1]{\texttt{DAC-#1}}
\newcommand{\umolJ}{\ensuremath{\mu}mol/J}

\title{Does the cold-start insight for direct air capture survive exergy accounting?\\
A negative result for our hypothesis, with a certified surrogate and a Lean core}
\author{Xavier Callens (SocrateAI Lab)}
\date{Preprint v1, 2026-10-27}
\hypersetup{pdftitle={Does the cold-start insight for direct air capture survive exergy accounting? A negative result for our hypothesis},
  pdfauthor={Xavier Callens}}

\begin{document}
\maketitle
\begin{center}\small Zenodo DOI to be assigned on release \quad (not peer reviewed)\end{center}

\noindent\textbf{AI-assistance disclosure.} This preprint, its Lean code, its Python code and its claim ledger were
produced by AI agents (Claude, Anthropic) under the author's direction. The author is responsible for the content. No human
has audited the Lean statements for adequacy, and no human referee has checked the hand arguments.
\textbf{No independent audit of any kind (human or AI) exists at the time of writing}; the producer-is-not-verifier recompile
is pending, and no statement lock exists. The source authors (NIST) have not been contacted and have not seen this work.

\begin{abstract}
McDannald et al.\ (arXiv 2501.04825) define an intrinsic capture efficiency $\xi=\Delta n_{\rm CO_2}/E_{\rm total}$ for a sorbent in an
idealised piston cycle and conclude that it is ``greatly beneficial to start cold'': 11 of the 14 Pareto rows of their Table~2 start
at $T_1\le 217.7$ K, and all start at or below $236.6$ K. Their $E_{\rm total}$ is a face-value sum of heats and works: the cooling of Step~1 is counted as a heat of adsorption at $T_1$,
the return Step~3 is set to zero (eq.~5), and no ambient temperature appears. We asked whether the conclusion
survives when every heat is weighted by its Carnot factor against an explicit ambient $T_0$, sub-ambient cooling is priced, and the
return step is charged. \textbf{Our hypothesis was that it would not (prior $0.35$); it is refuted.} For all 14 Table~2 rows the exergy-optimal start
temperature, found by a joint grid over $T_1$, $\Delta T$ and $P_{\rm end}$, lies between $200$ and $236.5$ K in every one of 18 variants
($T_0\in\{273.15,298.15,313.15\}$ K, three heat-recovery levels, with and without a purity constraint). The benefit of a cold start shrinks
(best $\xi$ with a cold start over best $\xi$ with $T_1\ge298.15$ K: face value $10$--$819$, exergy without recovery $3.5$--$129$,
with full reversible recovery $1.25$--$6.6$) but stays above $1$. Exergy-weighted efficiencies on the stored paths are $25.4$--$36.8$~\umolJ{}
against $5.98$--$11.26$ at face value (nominal recomputation; printed posterior means $6.12$--$11.259$). The published $\xi$ never
exceeds $1/W_{\min}$ of the ideal mixture; five dataset entries (CSMNOX, GUSNEN, PACYIE, SABWAU01, UJOKIP05) have heat flows that are
inconsistent with the second law, which we read as inconsistent inputs in those entries.

\textbf{Tier split.}
\tA{} (Lean~4.34.1, axioms within $\{$\lean{propext}, \lean{Classical.choice}, \lean{Quot.sound}$\}$): the twelve theorems of
\lean{DacExergy/Core.lean}, which are real-number algebra (Carnot weights in $(0,1)$, a factorisation, sign and threshold equivalences,
positivity and monotonicity of $RT_0\ln(1/x)$); the planted controls are caught by the axiom audit. No statement lock; no independent recompile yet.
\tL{}: the source paper's model (eqs.~1--10, Table~2) and the dataset and code licences, taken as read. \tC{}: the closed-form survival conditions
(an N$_2$ threshold and a heat-capacity threshold) and mpmath-interval enclosures, about a \emph{stylised two-component surrogate only}.
\tX{}: the full-model re-accounting, the grid scans, the sample, and the dataset-wide checks.
Every claim in the text carries one of these tags and a ledger identifier (\id{A-nnnn} etc.).

\medskip
\noindent\fbox{\parbox{0.96\linewidth}{\small\textbf{What this is NOT.} These are thermodynamic limits of a published \emph{simplified
model} (idealised piston, Henry/IAST loadings, binary CO$_2$/N$_2$ ``air'', no water, no heat recovery in the original). It is
\textbf{not} a statement about real plants, costs, deployment, sorbent performance in practice, or any climate effect, and it makes
\textbf{no recommendation}. It does not say the NIST metric or its authors are wrong: face value and exergy are two accounting conventions and the
source states its own scope. Nothing here has been discussed with the source authors.}}
\end{abstract}

\paragraph{Markers and ledger.} \tA{}: kernel-checked Lean, axioms within the allowlist, controls caught; the tag attaches to the Lean statement, never to its physical gloss.
\tL{}: a cited source taken as stated. \tC{}: executed interval code with a written argument, not machine-checked, \emph{about the surrogate}.
\tX{}: floating-point computation on the real model data; it supports nothing by itself beyond what it shows. Identifiers \id{A-nnnn}, \id{L-nnnn}, \id{C-nnnn}, \id{X-nnnn} refer to
\texttt{ledger/ledger.json} (Section~\ref{sec:ledger}).

%=====================================================================
\section{Setting, source and preregistration}\label{sec:setting}
This is item K1 of a computational programme (Proposal~4, preregistered 2026-10-27, \texttt{results/proposal4/2026-10-27/preregistration.json}).
The source is McDannald, Siderius, DeCost, Choudhary and Ortiz-Montalvo, \emph{Intrinsic Direct Air Capture}~\cite{mcd}
(arXiv:2501.04825v1, 8 Jan 2025; the only version on arXiv), with its data~\cite{zen} (Zenodo 10.5281/zenodo.14452152, from the NIST Data Publication
10.18434/mds2-3687, CC-BY-4.0) and its Julia code~\cite{nist} (NIST IntrinsicDACCycle v0.0.2, NIST software licence). We credit the authors and NIST fully:
every number below that is not labelled ours is theirs, and the headline of their paper is a thermodynamically careful model whose scope
(no heat recovery, intrinsic comparison of sorbents) they state themselves. We read the paper in full (28 pages) and
the relevant Julia sources line by line. \textbf{The supplement was not obtainable} (not on arXiv, not in the Zenodo record, not in the code repository) and was not read.

The preregistered item asked: (K1-H1) under exergy-consistent accounting the Pareto-optimal start temperature moves to $T_1\ge T_0$ for most materials (prior $0.35$);
(K1-H2) no dataset entry has exergy-weighted efficiency above $1/W_{\min}$, and any violation flags a model-input inconsistency (definite answer, prior $0.45$);
(K1c) the text value ``$1.0\times10^5$ mol/J'' versus Table~2 in \umolJ{}; (K1x) sensitivity to $T_0$ and heat recovery.
Kill rules: if the dataset and code lacked per-path temperature-resolved energy terms, H1 would stay a Tier~C argument; if the (unread) supplement showed that
sub-ambient cooling is already priced, H1 would be withdrawn. Neither fired: the dataset exposes per-path arrays, and in the \emph{code} the Step~3 term is $E_3=0$ in all four evaluators. The supplement being unread, we cannot say it is not
priced elsewhere in the paper's text; the eq.~5 text says it is zero.
Outcome: H1 refuted, H2 answered (the published $\xi$ never exceeds $1/W_{\min}$; five entries flagged).

%=====================================================================
\section{A step-by-step account for non-specialists}\label{sec:human}
\paragraph{Step 1: the cycle as published.}\tL{} A sorbent (here a metal-organic framework, MOF) sits in a piston. \emph{Step~1:} it is exposed to ``air'' (400 ppm CO$_2$, balance N$_2$) at
temperature $T_1$ and pressure $P_1$ and reaches equilibrium, releasing heat of adsorption that is carried away to keep $T_1$ fixed. \emph{Step~2:} the chamber is closed and the temperature is raised and the
pressure lowered along a path to $(T_{\rm end},P_{\rm end})$; gas desorbs and leaves through a check-valve. \emph{Step~3:} the system returns to $(T_1,P_1)$. The captured amounts
are $\Delta n_{\rm CO_2}$, $\Delta n_{\rm N_2}$ (mol/kg), the purity is $x_{\rm end}=\Delta n_{\rm CO_2}/(\Delta n_{\rm CO_2}+\Delta n_{\rm N_2})$ (eq.~7) and the metric is
\[\xi=\frac{\Delta n_{\rm CO_2}}{E_{\rm total}}\quad(\text{eq.~8}),\qquad E_{\rm total}=E_1+E_2+E_3\quad(\text{eq.~6}).\]

\paragraph{Step 2: the energy balance as read.}\tL{} Eq.~1: $E_1=Q_{\rm ads}-W_1$, which in the text becomes
$E_1=\Delta n_{\rm CO_2}\bigl(q_{\rm ads,CO_2}(T_1)+RT_1\bigr)+\Delta n_{\rm N_2}\bigl(q_{\rm ads,N_2}(T_1)+RT_1\bigr)$: a heat removed at $T_1$ plus a compression work.
Eq.~2: $E_2=Q_{\rm ads,CO_2}+Q_{\rm ads,N_2}-W_{2,{\rm CO_2}}-W_{2,{\rm N_2}}+E_{\rm heat,CO_2}+E_{\rm heat,N_2}+E_{\rm heat,sorb}+E_{\Delta P}$, and eq.~4 writes it as an integral over a path $s$ with
$dT/ds\ge0$, $dP/ds\le0$ (eq.~3), using heat capacities $\tfrac92R$ and $\tfrac72R$ for adsorbed CO$_2$ and N$_2$ and a fitted sorbent heat capacity $C_{V,{\rm sorb}}(T)$; (the argument of the logarithm in
the pressure term is not legible in our text extraction, so we read it in the Julia code and do not reproduce it here). Eq.~5: $E_3=0$ (``we assume no energy recovery''; the text adds that real systems would recover some). Eq.~9 is $C=1/\xi$.
Eq.~10 is the optimisation prior: $T_1\ge200$ K, $\Delta T\sim U(0,200)$, $P_1\sim U(101325,\,1.1\times101325)$ Pa, $\Delta P\sim U(-101225,0)$ Pa. Our Python
re-implementation follows the Julia functions line by line (\texttt{data/code/k1/dac\_model.py}); the Julia code is not redistributed here (Section~\ref{sec:licence}).

\paragraph{Step 3: why heats at face value differ from exergy.}\tL{}\,\tC{} A joule of heat is not a joule of work. Heat $Q$ delivered at temperature $T>T_0$ (with $T_0$ the ambient temperature) is worth at most
$Q(1-T_0/T)$ of work, the Carnot factor; this is the cost of producing that heat reversibly from work. Removing heat $Q$ \emph{below} ambient (e.g.\ cooling the sorbent to $T_1\approx200$ K while it adsorbs) needs a refrigerator, which costs at least $Q(T_0/T-1)$ of work;
heat rejected above ambient can be thrown away (no-credit accounting, $r=0$) or partly turned back into work (full reversible recovery, $r=1$). The paper's face-value sum treats all three alike
and sets the return step to zero, which is the point of our audit. We therefore (i) fix an explicit $T_0\in\{273.15,298.15,313.15\}$ K; (ii) record every heat once, with its direction and temperature (Step~1: $Q_{\rm ads}$ removed at $T_1$; Step~2: the per-step
heat at the mid-step temperature; Step~3: sensible cooling of the sorbent and of the adsorbed gas from $T_{\rm end}$ back to $T_1$, integrated with the same heat capacity $C_V(T)$ by an independent trapezoid, an \emph{explicit return step}); (iii) weight each by $\lvert1-T_0/T\rvert$ with the sign convention above; (iv) keep the works as published. The cost is
$B(r)=B_{\rm NC}-r\cdot({\rm credits})$ per mole of CO$_2$ and $\xi_{\rm ex}=1/B$. Lean states the two weights (\lean{carnot\_in\_mem\_unit}: $0<1-T_0/T<1$ for $T>T_0$, \id{A-0001}; \lean{refrig\_weight\_lt\_one}: $0<T_0/T-1<1$ for $T_0/2<T<T_0$) \tA; that these are the right weights is the textbook argument, not a Lean statement.

\paragraph{Step 4: the reversible floor.}\tA{}\,\tX{} Even a perfect machine must pay the ideal-mixture minimum work to extract CO$_2$ from the mixture. For a pure product at mole fraction $x$ this is
$W_{\min}=RT_0\ln(1/x)$, which is $19.4$ kJ/mol at $x=400$ ppm and $T_0=298.15$ K (hand arithmetic; $1/W_{\min}\approx51.6$~\umolJ{}). Lean proves that this function is positive for $0<x<1$ and strictly decreasing in $x$ (\lean{wmin\_pure\_pos}, \lean{wmin\_pure\_strictAnti}, \id{A-0004}); for less pure products our code uses
the corresponding mixture expression (Tier~X, not in Lean). A model whose reversible exergy cost $B_{\rm rev}$ (all credits taken) falls \emph{below} $W_{\min}$ would violate the second law, so it signals inconsistent inputs. This is H2.

\paragraph{Step 5: the closed-form thresholds in a stylised surrogate.}\tC{} To see \emph{why} a cold start can still win, we use a two-component surrogate of the cycle: van~'t~Hoff Henry loadings for CO$_2$ and N$_2$ with heats of adsorption $q_C,q_N$, a sensible-heat
term $c=\gamma Rn_0$ and a recovery fraction $r$ on the sensible heat, started at $T_c<T_0$ instead of $T_0$ and heated to the same $T_h$. The benefit is
\[\beta=C(T_0)-C(T_c)=A+\lambda B+(1-r)\,\gamma R\,D,\qquad A_{\rm ex}=-(T_0-T_c)\Bigl(\frac{q_C}{T_c}-R\Bigr),\quad A_{\rm fv}=R(T_0-T_c),\]
where $\lambda$ is the N$_2$-to-CO$_2$ loading ratio at the inlet and $B,D$ are differences of loading-ratio-weighted terms (\id{C-0001}; derivation by hand, algebra in Lean). Beside each step the Lean statement is:
\begin{itemize}\itemsep=1pt
\item the CO$_2$ term alone penalises cold under exergy: $A_{\rm ex}<0$ whenever $q_C>RT_c$ (\lean{numerator\_factor}, \lean{co2\_term\_penalises\_cold}, \id{A-0002}); at face value it favours cold ($A_{\rm fv}>0$, \lean{co2\_term\_facevalue\_favours\_cold}). This is \emph{not} the sign of the whole benefit;
\item $\beta$ is affine in $r$ (\lean{benefit\_affine\_in\_r}), and when $\gamma D\ge0$ the case $r=1$ is the worst case for a cold start (\lean{worst\_case\_full\_recovery}, \id{A-0003});
\item \textbf{N$_2$ threshold} (at $r=1$): for $B>0$, cold wins iff $\lambda>\lambda^\ast=-A_{\rm ex}/B=(T_0-T_c)(q_C/T_c-R)/B$ (\lean{lambda\_threshold\_iff});
\item \textbf{heat-capacity threshold} (for $r<1$): with $N=-(A+\lambda B)$ and $RD$ replaced by $(1-r)RD>0$, cold wins iff $\gamma>N/((1-r)RD)$ (\lean{threshold\_iff});
\item a rational instance ($R=8.314$, $q=40$ kJ/mol, $T_c=200$ K, $T_0=298.15$ K) with negative CO$_2$-only term, checked by \lean{instance\_r1} (this is \lean{norm\_num}, an identity check and not a theorem about the world).
\end{itemize}
Reading: under exergy accounting the refrigeration of the CO$_2$ adsorption heat is a real penalty, and a cold start wins only because lowering $T$ changes the \emph{ratio} of N$_2$ to CO$_2$ uptake (the N$_2$ term $\lambda B$ and the sensible term). This is the paper's own mechanism (CO$_2$ uptake diverges from N$_2$ at low $T$); the purity at the optimum is $0.98$ at 200 K but $0.20$ at 240 K for LETQAE01 on its slice.

%=====================================================================
\section{Results}\label{sec:results}
\subsection{Positive control: Table 2 reproduced before any re-weighting}
\tX{} (\id{X-0001}) The Python re-implementation, fed the public dataset, reproduces the 14 rows of Table~2 in all eight printed fields each (14/14 rows, from the stored posterior means and standard deviations, which round to the printed strings), and recomputes every stored nominal array
(loadings, heats, every Step-1/Step-2 array, $E_1$, $E_2$, $E_{\rm total}$) with a maximum normalised deviation of $3.8\times10^{-11}$ (one number per material, the density, is back-solved from one stored Henry constant because the CSD material files are not in the public record, see Section~\ref{sec:dev}). The posterior statistics cannot be regenerated bit for bit (unseeded draws, Section~\ref{sec:obs}); the nominal $\xi$
differs from the posterior mean by $-2.3\%$ to $+1.3\%$. The Gaussian-process heat capacity re-fit agrees with the stored one to $4.0\times10^{-6}$ relative. A second, vectorised implementation (\texttt{fast\_costs.py}) agrees with the ledger evaluator to $2.3\times10^{-15}$.

\subsection{Exergy re-accounting of the 14 published paths}
\begin{table}[h]\centering\footnotesize
\caption{The 14 Table~2 rows (\tX; \id{X-0002}, \id{X-0003}, \id{X-0004}). $\xi_{\rm pub}$: printed posterior mean (Table~2). $\xi_{\rm fv}$: our nominal recomputation at face value. $\xi_{\rm NC}$ ($r=0$) and $\xi_{r=1}$: exergy-weighted $\xi$ on the \emph{stored} path at $T_0=298.15$ K, in \umolJ{}. Refrig.: refrigeration share of the no-credit exergy cost. $T_1^\ast$: exergy-optimal start ($r=0$, $T_0=298.15$ K, unconstrained joint grid). Last three columns: cold-start benefit ratio (best $\xi$ cold over best $\xi$ with $T_1\ge298.15$ K).}
\begin{tabular}{llrrrrrrrr}\toprule
material & $\xi_{\rm pub}$ & $\xi_{\rm fv}$ & $\xi_{\rm NC}$ & $\xi_{r=1}$ & refrig. & $T_1^\ast$ (K) & fv & $r{=}0$ & $r{=}1$\\\midrule
VANNIK & 11.259(38) & 11.26 & 36.5 & 38.9 & 0.69 & 200.0 & 30.6 & 11.9 & 5.3\\
VANNIK & 11.223(41) & 11.24 & 36.6 & 38.9 & 0.69 & 200.0 & 29.0 & 11.5 & 5.3\\
LETQAE01 & 11.22(18) & 11.21 & 36.8 & 38.5 & 0.76 & 202.0 & 16.4 & 10.1 & 6.6\\
BEVQID & 9.186(66) & 9.15 & 34.3 & 39.1 & 0.67 & 210.4 & 78.2 & 16.8 & 1.5\\
BEVQID & 9.156(68) & 9.12 & 34.8 & 39.2 & 0.66 & 210.1 & 82.8 & 18.2 & 1.5\\
FASJAL & 8.186(69) & 8.19 & 35.0 & 37.0 & 0.47 & 236.5 & 10.2 & 3.5 & 1.2\\
FASJAL & 8.148(62) & 8.14 & 35.1 & 37.7 & 0.48 & 235.0 & 12.2 & 4.1 & 1.3\\
FASJAL & 8.028(69) & 8.03 & 34.5 & 37.6 & 0.47 & 234.4 & 13.2 & 4.5 & 1.3\\
IFUDAO & 7.26(29) & 7.27 & 31.3 & 38.5 & 0.68 & 217.4 & 121.6 & 24.6 & 1.3\\
IFUDAO & 7.00(28) & 7.03 & 31.5 & 38.9 & 0.64 & 217.6 & 117.6 & 23.4 & 1.3\\
MAXHEA & 6.84(36) & 6.77 & 28.2 & 38.6 & 0.72 & 213.0 & 209.5 & 35.9 & 1.4\\
MAXHOK & 6.56(41) & 6.64 & 28.4 & 39.1 & 0.74 & 212.9 & 289.7 & 55.8 & 1.4\\
MAXHOK & 6.26(52) & 6.26 & 29.0 & 39.2 & 0.71 & 213.3 & 254.7 & 47.0 & 1.3\\
PARFOF & 6.12(52) & 5.98 & 25.4 & 40.0 & 0.80 & 203.4 & 819.2 & 129.3 & 1.5\\\bottomrule
\end{tabular}
\end{table}
Three readings, each \tX{} and each inside the published model. (1) \emph{Exergy-weighted efficiencies are higher than face value}, not lower: $\xi_{\rm NC}=25.4$--$36.8$~\umolJ{} against $\xi_{\rm fv}=5.98$--$11.26$ (nominal; printed $6.12$--$11.259$). A plausible reading (not separately tested) is that in this model the face-value sum is dominated by the sensible heating of the sorbent (the source says so, p.~17), and Carnot weights below 1 shrink that term by more than refrigeration adds. The face-value $\xi$ is an accounting-convention quantity and is \emph{not} an upper bound on CO$_2$ per unit of work; the source says as much (pp.~14--15: heat recovery ``could therefore exceed our stated upper limit''). (2) The reversible floor on the stored paths is
$B_{\rm rev}=25.0$--$27.0$ kJ/mol against $W_{\min}(P_0,y)=18.8$--$19.4$ kJ/mol, i.e.\ $\xi_{\rm rev}=37.0$--$40.0$~\umolJ{}, below $1/W_{\min}$ in all 14 rows. (3) Refrigeration is $47$--$80\%$ of the no-credit exergy cost: the penalty the hypothesis relied on is real and large.

\subsection{H1: the optimal start temperature does not move up}
\tX{} (\id{X-0003}) Joint re-optimisation (Section~\ref{sec:dev} for the protocol and its amendment) over $T_1\in[200,320]$ K, $\Delta T\ge0$ with $T_{\rm end}\le400$ K and $P_{\rm end}\in[100,60000]$ Pa (the optimiser's 100 Pa bound), $P_1$ as stored, with refinement around the best cell of each accounting. In \textbf{14 of 14 rows} the exergy-optimal $T_1^\ast$ lies in $[200,\,236.54]$ K, for every
$T_0\in\{273.15,298.15,313.15\}$ K, every $r\in\{0,0.5,1\}$, with and without the purity constraint (18 variants); $0$ of $14$ rows have $T_1^\ast\ge T_0$ in any variant, against the preregistered criterion ($>50\%$ at $T_0=298.15$ K, $r=0$). The face-value control puts all 14 optima below 273 K (maximum $231.98$ K), as it must for the evaluator to be trusted.
What \emph{does} move is the hot end: unconstrained optimal $T_{\rm end}$ is $227$--$400$ K at face value and $301$--$340$ K under no-credit exergy accounting (\id{X-0004}), because heating near ambient is cheap in exergy but a refrigerated start is not avoided. The warm alternative for comparison is the best path with $T_1\ge298.15$ K: its purity collapses (for VANNIK, $5\%$ at $T_1=298$ K), consistent with its losing even when refrigeration is charged.

\subsection{The cold-start benefit shrinks but stays}
\tX{} (\id{X-0004}) In the last three columns of Table~1 the ratio falls from $10$--$819$ at face value to $3.5$--$129$ (exergy, $r=0$), $2.6$--$80$ ($r=0.5$) and $1.25$--$6.6$ ($r=1$). At full reversible recovery several rows approach $1$ (FASJAL $1.2$--$1.3$); this is the regime where an exergy audit could most plausibly overturn the conclusion, and it does not here, within this model and grid. We did not test any process-level recovery.

\subsection{H2: no published \texorpdfstring{$\xi$}{xi} above \texorpdfstring{$1/W_{\min}$}{1/Wmin}; five inconsistent entries}
\tX{} (\id{X-0005}) Over the 8,765 materials with valid paths (15,161 paths; the public record holds 11,656 optimised files) the largest published face-value $\xi$ is $11.259$~\umolJ{} (VANNIK), against $1/W_{\min}\approx51.6$~\umolJ{}: no violation of the preregistered form.
The reversible exergy floor falls below the ideal-mixture $W_{\min}$ (lowest-pressure bound) for 18 path/$T_0$ pairs at $T_0=298.15$ K (14 at $273.15$ K, 21 at $313.15$ K), in five materials at $298.15$ K: CSMNOX, GUSNEN, PACYIE, SABWAU01, UJOKIP05. Most are degenerate paths (purity below $4\%$, $\xi\approx0$). SABWAU01 (purity $0.91$, $B_{\rm rev}=-42$ J/mol) and UJOKIP05 (purity $0.67$) are not.
We present these as \emph{inconsistent inputs in those entries} (a suggested cause, not tested, is extrapolated heats of adsorption); we do not claim that the method or the source model is at fault. Our gate rejects exactly these pairs (G3). A counting note: the first run script skipped cases with $W_{\min}(\text{low }P)<0$ in its own counter and reported 2 and 5 low-pressure violations at $298.15$ and $313.15$ K; the violator list (\texttt{h2\_violators.json}) counts them all and is the authoritative one.

\subsection{Beyond the 14 rows: a seeded sample and the whole record}
\tX{} (\id{X-0006}) A slice scan (T$_{\rm end}$, $P_1$, $P_{\rm end}$ held at the stored values) on a seeded random sample of 150 materials (seed $20261027$; each represented by its stored Pareto path of highest face-value $\xi$) finds $T_1^\ast\ge T_0$ in only \textbf{3 of 150} at $T_0=298.15$ K, $r=0$. Re-selecting among all 15,161 stored Pareto paths, the no-credit-exergy-best path starts below $298.15$ K for $8{,}420$ of $8{,}765$ materials (face value: $8{,}416$). Both are conditional on the stored paths and are not re-optimisations; the sample is not a population estimate.

\subsection{The surrogate: where it agrees with the model and where it does not}
\tC{} (\id{C-0002}--\id{C-0004}) Over the box $T_c\in[200,272]$, $T_0\in[273,313]$, $T_h\in[315,400]$ K, $q_C\in[12,67]$ kJ/mol, $A_{\rm ex}<0$ (certified; trivial since $q_C\ge12$ kJ/mol exceeds $RT_c\le2.3$ kJ/mol). At the surrogate parameters extracted from 164 parameter sets (the 14 paths plus the sample), the certified signs of $\beta$
at $T_0=298.15$ K, $r=0$ are $151$ positive (cold wins), $8$ negative, $0$ undecided (159 sets with $T_c<T_0$); at $r=1$, $131$ positive and $28$ negative. Intervals are outward-rounded mpmath ($80$ bits; python-flint is not installed). An independent plain-float implementation lies inside the enclosures for 2854 of 2856 points; the other two differ by $3\times10^{-8}$ relative from
float cancellation, with the same sign. (The first surrogate run omitted the factor $R$ on the sensible term; this float check caught it with 1904 of 2856 mismatches and the run was repeated.)

\tX{} (\id{X-0007}) The \emph{link} between surrogate and model is empirical: the surrogate sign agrees with the full-model slice scan on $149/150$ ($r=0$), $146/150$ ($r=0.5$) and $128/150$ ($r=1$) of the sample at $T_0=298.15$ K. On the 14 Table~2 rows the surrogate is \textbf{wrong at $r=1$ for most rows}: it disagrees on $11$ of $14$ at $T_0=273.15$ K, $10$ of $13$ scorable
rows at $298.15$ K and $5$ of $8$ at $313.15$ K (it predicts that cold loses, the full model says cold wins); at $r=0$ all scorable rows agree and at $r=0.5$ one material (FASJAL) disagrees. (The count falls because rows with no grid point on one side of $T_0$ are not scorable; the figures come from a per-row recount, whereas the run's own agreement file uses the first path of each material.) The certified statements are therefore about the surrogate only, and the surrogate is a fair proxy of the published model only for $r\le0.5$.

%=====================================================================
\section{The headline, and its tiers}\label{sec:head}
\begin{center}\footnotesize\begin{tabular}{p{7.2cm}p{0.9cm}p{5.0cm}}\toprule
statement & tier & basis\\\midrule
the source model, eqs.~1--10, Table~2, licences & \tL & arXiv:2501.04825v1; Zenodo; NIST (\id{L-0001}, \id{L-0002})\\
Carnot weights in $(0,1)$; $A_{\rm ex}<0$ for $q>RT_c$; affine in $r$; $\lambda$ and $\gamma$ threshold equivalences; $RT_0\ln(1/x)$ positive, decreasing & \tA & \lean{DacExergy/Core.lean}, 12 theorems (\id{A-0001..0004})\\
decomposition $\beta=A+\lambda B+(1-r)\gamma RD$ and its thresholds in the surrogate & \tC & hand derivation $+$ Lean (\id{C-0001})\\
certified surrogate signs at 164 parameter sets; float cross-check & \tC & mpmath intervals (\id{C-0003}, \id{C-0004})\\
Table~2 reproduced; exergy re-accounting; joint optimum $T_1^\ast\in[200,236.5]$ K; benefit ratios; H2; sample; surrogate agreement & \tX & grids and scans (\id{X-0001..0007})\\
\textbf{For the 14 published paths, in the published simplified model, the cold-start conclusion survives exergy accounting; our hypothesis H1 is refuted} & \tX & \id{X-0011}, resting on \id{X-0001}, \id{X-0003}, \id{X-0004}, \id{X-0008}\\\bottomrule
\end{tabular}\end{center}
The headline is filed Tier~X because it rests on grid optimisation of the real model, not on a proof; the ledger gate would block it if it were filed above any claim it depends on.

%=====================================================================
\section{Controls}\label{sec:controls}
All controls were run before the exergy numbers were interpreted (\id{X-0008}, \id{X-0001}, \id{A-0005}).
\paragraph{Positive controls (all met).}
\begin{itemize}\itemsep=0pt
\item \textbf{Preregistered: Table~2 reproduced to the printed digits} before any re-weighting: 14/14 rows, 8/8 fields (Section~3.1).
\item \textbf{Preregistered: an all-$T\ge T_0$ path.} A path entirely at $T\ge T_0$ (300 K start) gives heat-in weights all in $(0,1)$, no refrigeration term, and heat-out free. This matches \lean{carnot\_in\_mem\_unit} \tA.
\item The real BEVQID ledger passes the consistency gate (G1--G4).
\item The face-value joint optimum reproduces ``start cold'' (14/14 below 273 K); a stored path re-evaluated from raw inputs with the re-fitted heat capacity agrees to $4.2\times10^{-7}$ relative.
\end{itemize}
\paragraph{Negative controls (all four rejected by the gate, as preregistered).}
\begin{itemize}\itemsep=0pt
\item \textbf{Double-counted sorbent heating}: rejected by G2 (the sensible closure between Step~2 and Step~3 is off by $100\%$).
\item \textbf{Planted $\times10^{6}$ unit error}: rejected by G4 (reported $9.1\times10^{6}$ against the ledger $9.15$~\umolJ{}, and above $1/W_{\min}=51.8$~\umolJ{}).
\item \textbf{The sentence's value taken literally} ($1.0\times10^5$ mol/J, i.e.\ $10^{11}$~\umolJ{}): rejected by G4.
\item \textbf{Heats of adsorption removed} (a second-law violation planted): rejected by G3 ($B_{\rm rev}=8.7$ kJ/mol $<W_{\min}=13.5$ kJ/mol).
\item Lean planted controls (\id{A-0005}): \lean{ctl\_sorry} and \lean{ctl\_axiom} compile with rc~0 and are caught by \lean{\#print axioms} (\lean{sorryAx}; \lean{DacExergy.ctl\_cheat}).
\end{itemize}
\paragraph{The gate} (G1--G4) is ours: G1 every heat flow recorded once, G2 sensible-heat closure between Step~2 and Step~3, G3 $B_{\rm rev}\ge W_{\min}$ at the delivered purity, G4 reported $\xi$ equals $\Delta n/E_{\rm total}$ in the stated units and lies below $1/W_{\min}$. It was applied to all 14 Table~2 rows (14 pass at $T_0=298.15$ K) and to the whole record in H2.

\paragraph{The unit statement, stated neutrally.}\tL{}\,\tX{} (\id{L-0003}, \id{X-0009}) On p.~20 the text reads ``a few sorbents with intrinsic capture efficiencies above $1.0\times10^{5}$ mol/J''. Table~2 gives $\xi$ in \umolJ{} (largest $11.259$); the dataset field \texttt{Intrinsic\_capture\_efficiency} carries the units string ``mol/J'' and its largest posterior value is $1.1259\times10^{-5}$ mol/J; 60 of the 8,765 materials have a path above $1.0\times10^{-5}$ mol/J.
The next sentence compares with $2.46\times10^{-6}$ mol/J (US grid), which is meaningful only at the $10^{-5}$ scale, and the exponent of that number prints with its minus sign in the same text extraction. The printed sentence is therefore consistent with $1.0\times10^{-5}$ mol/J; it reads as a missing minus sign in a typeset exponent, whereas the table, the data and the later sentences are self-consistent. We checked this against the dataset's units string and two text extractions of the PDF, \emph{not} against the supplement (not obtained) and not against the page image. We drew no conclusion from it other than the negative control above; it is not an error of the method and no author has been contacted.

%=====================================================================
\section{Protocol, deviations and observations}\label{sec:dev}
\paragraph{Protocol amendment, written before the joint scan.} The operational test (\texttt{logs/h1\_operational\_test\_written\_before\_run.md}, written before any exergy number) scanned $T_1$ on a 1 K grid from 200 K while holding $(T_{\rm end},P_1,P_{\rm end})$ at the stored values. After that slice scan, and before the joint scan, the protocol was amended (\texttt{logs/h1\_protocol\_amendment\_2026-10-27T1300Z.md}): the slice measures the optimum along a slice that NSGA-III had chosen for a cold start, so warm starts with deeper vacuum and pure vacuum swings at $T_1=T_{\rm end}\ge T_0$ had never been evaluated.
The amended score is the joint grid (coarse: $T_1\in\{200,\ldots,320\}$ K step 10, $\Delta T$ step 10 K with $T_{\rm end}\le400$ K, $P_{\rm end}\in\{100,10^3,5\times10^3,2\times10^4,6\times10^4\}$ Pa; refinement around the best cell of each accounting). The slice result is kept and reported as conditional (Section~3.6). We report both results, and the criterion ($>50\%$ at $T_0=298.15$ K, $r=0$) was not changed.
\paragraph{Deviations from the preregistration} (the preregistered text of K1 is unedited).
\begin{enumerate}\itemsep=0pt
\item \textbf{Material files missing.} The CSD material files (\texttt{cell\_density}, \texttt{cell\_mass}) are not in the public record. The density is back-solved from one stored Henry constant per material (every other array is an independent check); scans use the cell mass $\rho V_{\rm box}$ from the saturation box for the saturation truncation. The location of the saturation edge (e.g.\ BEVQID's effective $T_1\approx209$--$210$ K) depends on this substitute; the direction of the H1 result does not.
\item \textbf{Table~2 reproduced from stored values}, plus an independent nominal recompute; the posterior itself is not regenerable (unseeded draws).
\item \textbf{Search space}: a 1-D slice, then a joint grid with refinement instead of NSGA-III over 8,759 materials (out of budget on 2 cores). The 150-material sample and the stored-Pareto re-selection are conditional. The first joint run let the refinement go to $P_{\rm end}=50$ Pa, below the optimiser's 100 Pa bound; it was rerun with the bound enforced, conclusions identical (the first run is kept as \texttt{scan\_joint\_v1\_Pend\_unbounded}).
\item \textbf{No python-flint}: mpmath intervals instead. \textbf{Lean runner}: the pinned wrapper \texttt{lean\_pinned.py}, not \texttt{anse/formal/lean\_runner.py} (the repository's usual route); the rc and the axiom output were checked. No leanchecker replay of \texttt{DacExergy} was run (the sibling C1 paper had one).
\item \textbf{Supplement not read} (K1c required checking the sentence against it before any mention; we mention it anyway, as instructed by the author, on the evidence above). No author contact (the proposal reserves that decision to the owner).
\item \textbf{GATE-1} (the statement-adequacy and vacuity checks added by the proposal addendum) was not run. Several Lean theorems are one-line tactic identities (\lean{ring}, \lean{norm\_num}); \lean{instance\_r1} and \lean{benefit\_affine\_in\_r} should be read as such, and a linter notes an unused hypothesis (\lean{h0}) in \lean{refrig\_weight\_lt\_one}.
\item \textbf{Counts}: the preregistration assumed one 11,660-material set; the public record holds 11,656 optimised files, of which 8,765 have valid paths. The \texttt{h2\_all.py} own-counter under-counted (Section~3.5).
\item \textbf{Shell discipline}: the author's isolated-shell rules (no variables, loops or heredocs) were lapsed in a few commands while preparing this paper (heredocs for several read-only inspection commands, one compound \texttt{cd}, one command with a shell variable that the guard refused); none touched git.
\end{enumerate}
\paragraph{Neutral observations on the published code and data}\label{sec:obs}\tX{} (\id{X-0010}); none is an error of the method, and none is a claim about the authors' intent.
\begin{itemize}\itemsep=0pt
\item The posterior function draws 100 samples without a seed and, inside the loop, reuses one heat-capacity draw, so printed $\xi$ uncertainties omit $C_V$ variability and the means are not bit-reproducible; in the optimised JSON files \texttt{Captured\_CO2} is the posterior mean while \texttt{E\_Balance} is nominal, so they should not be divided by each other.
\item Three descriptions of the single path of Fig.~4/Table~1: the text, 250$\to$350 K then vacuum to 40,530 Pa; the stored SinglePath files and the dataset README, 250$\to$375 K and 42,844.84 Pa (checked on BEVQID, LETQAE01, VANNIK); the code default, 300$\to$350 K.
\item Eight Table~2 rows have $P_1$ below 101,325 Pa (82,466--101,126 Pa) although eq.~10 and the v0.0.2 optimiser bound $P_1$ to between 101,325 and 111,457.5 Pa; the stored results were probably produced with other bounds or another code version.
\item With the published evaluator, our joint grid finds nominal paths that exceed several Table~2 rows in both $\xi$ and purity (PARFOF $8.42$ against $5.98$~\umolJ{} at purity $0.99999$; MAXHEA $8.62$ against $6.77$; IFUDAO $8.60$ against $7.27$; nominal against nominal). The published front \emph{may not be fully converged}; this is an observation to check, not a claim. Those face-value optima end below ambient ($T_{\rm end}$ 235--285 K), so their heating is cheap in exergy.
\end{itemize}

%=====================================================================
\section{Lean formalisation}\label{sec:lean}
Modules \path{formal_cert/Programmes/DacExergy/} (copies in \texttt{data/lean/}): \texttt{Core.lean} (12 theorems, imports \lean{Mathlib.Tactic} and \lean{Mathlib.Analysis.SpecialFunctions.Log.Basic}; never \lean{import Mathlib}) and \texttt{Controls.lean} (never imported). Compiled through the pinned wrapper (Lean~4.34.1), rc~0; every theorem's \lean{\#print axioms} reports exactly
$[$\lean{propext}, \lean{Classical.choice}, \lean{Quot.sound}$]$ (\texttt{data/logs/lean\_core.json}). No \lean{sorry} or \lean{native\_decide}. The planted controls compile with rc~0 and the audit exposes \lean{sorryAx} and \lean{DacExergy.ctl\_cheat} (\texttt{lean\_controls.json}).
sha256: \texttt{Core.lean} \texttt{4bfa4f7ebd979d53e8e98847369610339f5bc3530dbd7f43f983fbbfe1a6833c}; \texttt{Controls.lean} \texttt{25017505f1fd76b08024326b9debd03664179602272e025d1d4d1d4f0448aac8}.
\textbf{What Tier~A certifies:} twelve statements about real numbers. \textbf{It does not certify} that the weights are the right physics, that the surrogate decomposition matches the published model, any grid optimum, or any number in Section~3; the transcendental parts (exponentials in the van~'t~Hoff loadings, logarithms in the sensible integral) are interval arithmetic outside the kernel (Tier~C). Note that \lean{co2\_term\_penalises\_cold} is explicitly \emph{not} the sign of the whole benefit.

%=====================================================================
\section{Claim ledger}\label{sec:ledger}
\texttt{ledger/build\_ledger.py} files $24$ claims: $5$ at Tier~A (\id{A-0001..0005}), $3$ at L (\id{L-0001..0003}), $4$ at C (\id{C-0001..0004}) and $12$ at X (\id{X-0001..0012}). The Tier~C claims are phrased as statements about the surrogate at stated numeric parameter values, so they do not depend on the Tier~X claims that describe how the values were extracted from the model; the link is the separate Tier~X claim \id{X-0007}.
Each claim has a content-addressed evidence blob in \texttt{ledger/evidence/}; the headline is \id{X-0011}; \id{X-0012} records the absences. The gate (\texttt{tools/ledger.py ledger.json -{}-evidence-dir evidence}, verbatim in \texttt{ledger/gate\_output.txt}) exits with code~\textbf{1}: \textbf{0 blocking findings and 5 flags} (\texttt{LEDGER\_UNAUDITED\_TIER\_A}, one for each Tier~A row: no \emph{person}
has certified that the Lean statements mean what their glosses say). Exit code~1 is the honest state, not a failure to be cleared by editing. There is no AI audit either, and no statement lock.

%=====================================================================
\section{Limitations}\label{sec:lim}
\begin{itemize}\itemsep=1pt
\item \textbf{Everything is inside a published simplified model}: Henry/IAST loadings, absolute adsorption, adsorbed-gas heat capacities $\tfrac92R$ and $\tfrac72R$, GP-extrapolated heat capacity, binary CO$_2$/N$_2$ air, no water, ideal piston. A different model could give a different answer; nothing transfers to real plants.
\item \textbf{Exergy weighting uses ideal reversible machines against one ambient $T_0$}; works are kept at face value; no process-level recovery was modelled; $r$ is a bookend, not a design.
\item \textbf{The joint search is a grid with refinement, not a proof of global optimality}; the cell-mass substitute shifts the saturation edge.
\item \textbf{The interval results are about the surrogate} and the surrogate is wrong on 11 of 14 Table~2 rows at $r=1$ and $T_0=273.15$ K (Section~3.7); the Lean results certify algebra only.
\item \textbf{The supplement was not obtained and not read}; the sub-ambient pricing question is settled only for the code and the main text.
\item \textbf{No independent audit}: no human or AI auditor has reviewed this paper, the code or the Lean statements; the producer $\neq$ verifier recompile is pending; no statement lock; GATE-1 not run. The advisor review of the K1 lane timed out and produced no output.
\item \textbf{Repository guards were not run} (\texttt{antigravity\_guard.py}, \texttt{test\_rigor\_guard.py}, \texttt{pytest tests/}); no guard output is claimed.
\item \textbf{No author contact} and no statement about the authors' intent. Wording that implies the NIST bound is wrong is not supported and is not intended; the metric is a face-value convention by the source's own description.
\item \textbf{Novelty is unchecked} beyond the source paper read; this is not a priority claim.
\item \textbf{Data note}: the run note quotes ``11/14'' surrogate disagreement without a temperature; the per-row recount shows it holds at $T_0=273.15$ K, with $10/13$ at $298.15$ K.
\end{itemize}

%=====================================================================
\section{Credit, licences and reproducibility}\label{sec:licence}
\paragraph{Credit.} The model, the code, the dataset and Table~2 are the work of A.~McDannald, D.~W.~Siderius, B.~DeCost, K.~Choudhary and D.~L.~Ortiz-Montalvo at the National Institute of Standards and Technology. We acknowledge NIST as the source of the data and software.
\paragraph{What is and is not redistributed.} \emph{Not redistributed}: the NIST dataset (3.4~GB; CC-BY-4.0, available at Zenodo~\cite{zen}); the NIST Julia code; the paper PDF or its text. The 51 dataset files used for the Table~2 materials are listed by name and sha256 in \texttt{data/zenodo\_members\_SHA256SUMS.json} (they are available from the public record, and we keep this bundle small and free of third-party data);
the full archive has sha256 \texttt{75a73c75a376a4611511af4a9e99b15af10096fa1886e5d969e4a1d8268df3fa} and md5 \texttt{a2520f520140c60f8a4663d5b403966c}, the Zenodo value. \emph{Redistributed with this paper}: our own Python code (\texttt{data/code/k1/}), which re-implements the NIST Julia functions in Python.
This code is a derivative of NIST-developed software: it was \textbf{changed} (translated from Julia to Python, with additional exergy accounting) by the author's AI agents on 2026-10-27, and the NIST licence notice is kept in \texttt{data/code/k1/NIST\_NOTICE.txt}; the software is provided ``as is'' and neither NIST nor the authors warrant its correctness. We also distribute small logs, the Lean sources, and the run note.
\paragraph{Reproducibility.} Repository AutoevolveAI, branch \texttt{worktree-openai-math-discovery}; K1 lane committed at \texttt{e298698}, HEAD when this paper was written \texttt{e29869853ec8056f6b8aaa1c32e09e018b80ae61}; this paper's files are uncommitted working-tree files and nothing was pushed or published. The paper's \texttt{data/} holds the code, Lean sources, small logs, and the large logs \texttt{cert\_sign.json} and
\texttt{scan\_T1.json}; the 3 MB \texttt{h2\_all.json} and \texttt{scan\_joint\_v1\_Pend\_unbounded.json} stay in \texttt{results/proposal4/2026-10-27/K1/logs/} (sha256 \texttt{a3a59822\ldots} and \texttt{3e615bf5\ldots}). Commands (from the repository root):
\begin{lstlisting}[basicstyle=\ttfamily\scriptsize]
python -I code/run_pc.py ; python -I code/run_exergy_table2.py          # positive control, exergy re-accounting
python -I code/run_scan.py ; python -I code/run_joint.py                # slice scan, joint grid (amended)
python -I code/h2_all.py ; python -I code/h2_violators.py               # H2
python -I code/surrogate_params.py ; python -I code/cert_sign.py        # surrogate, interval signs
python3 scripts/certified_numerics/lean_pinned.py compile formal_cert/Programmes/DacExergy/Core.lean
python -I data/code/table2_surrogate_check.py ... ; python -I data/code/recount_ranges.py ...   # paper-side recounts
python -I ledger/build_ledger.py ; python -I ledger/run_gate.py         # gate: exit 1 = 0 blocking, 5 flags
\end{lstlisting}
The exact argument lists are in the script headers and the run note; the scans were not re-run for this paper (the paper-side recounts were). The run used two CPU cores and no GPU.

%=====================================================================
\begin{thebibliography}{99}\small
\bibitem{mcd} A.~McDannald, D.~W.~Siderius, B.~DeCost, K.~Choudhary, D.~L.~Ortiz-Montalvo, \emph{Intrinsic Direct Air Capture}, arXiv:2501.04825v1 (8 Jan 2025), National Institute of Standards and Technology. Source of the model, eqs.~1--10 and Table~2 (\tL). Supplementary information not read.
\bibitem{zen} A.~McDannald, D.~W.~Siderius, \emph{Data for Intrinsic DAC calculations}, version 1.0, Zenodo, doi:10.5281/zenodo.14452152 (NIST Data Publication doi:10.18434/mds2-3687), CC-BY-4.0.
\bibitem{nist} NIST, \emph{IntrinsicDACCycle} v0.0.2 (Julia package), \url{https://github.com/usnistgov/IntrinsicDACCycle}, NIST software licensing statement; read, not redistributed.
\bibitem{mathlib} The mathlib Community, \emph{The Lean mathematical library}, arXiv:1910.09336 (2019); cited for the library only.
\end{thebibliography}
\end{document}
